\documentclass{article}
\usepackage{amsmath}
\usepackage[a4paper]{geometry}
\usepackage{fancyhdr}
\pagestyle{fancy}
\lhead{Das Zerfallsgesetz}
\rhead{März 2026} 
\newcommand{\dd}{\text{d}} 
\begin{document} 
\section{Das Zerfallsgesetz}
Ist ein Anfangsbestand an Kernen $N_0$ oder die Anfangsaktivität $A_0$ und die Zerfallskonstante $\lambda$ bekannt, so kann mit dem Zerfallsgesetz die Kernenzahl beziehungsweise die Aktivität in der Zukunft, oder wenn es sachkontextlich Sinn ergibt, auch in der Vergangenheit, bestimmt werden. Es gelten
\[
 N(t) = N_0 \cdot e^{-\lambda \cdot t} 
 \qquad \text{und}
 A(t) = A_0 \cdot e^{-\lambda \cdot t}  
\] 
\subsection{Herleitung} 
Werden die beiden Definition der Aktivität gleichgesetzt, so folgt 
\[
 \frac{\dd N}{\dd t} = -\lambda \cdot N 
\]
Wird auf beiden Seiten durch $N$ dividiert und ein Integral genommen, so folgt
\begin{align*} 
 \int \frac{\dd N}{N} &= \int -\lambda \cdot dt \\ 
 \ln{N} &= -\lambda \cdot t + c \\
 N &= e^{-\lambda \cdot t} \cdot e^c
\end{align*}
Wird dies nun bei $t=0$ betrachtet, so ist $e^{-\lambda \cdot 0} = 1$, so dass $e^c=N_0$, die Anzahl der unzerfallenen Kernen am Anfang, sein muss.
 
Daraus folgt die Formel für $N(t)$, welches die Anzahl der noch vorhandenen, unzerfallenen, Kerne zum Zeitpunkt $t$ angibt.
\[
 N(t) = N_0 \cdot e^{-\lambda \cdot t} 
\]
Wird beides durch $t$ dividiert folgt das gleiche 
\[
 A(t) = A_0 \cdot e^{-\lambda \cdot t} 
\]
 
\subsection{Grenzen} 
Es muss angemerkt werden, dass das Zerfallsgesetz nur einen stochastischen Prozess darstellt. Dem \emph{Gesetz der großen Zahlen} und der riesigen Menge an Kernen in auch geringen Mengen an Stoff ist der Zerfallsgesetz sehr akkurat.
 
Ist in einem Beispiel der Anfangsbestant aber gewollt klein angegeben, so verliert das Zerfallsgesetz an genauigkeit.  
\end{document}